\section{Introduction}
Multimodal Large Language Models (LVLMs) have achieved remarkable progress in image understanding by leveraging the world knowledge of Large Language Models (LLMs). As a result, LVLMs can be applied to a wider range of vision-language tasks with stronger reasoning capabilities.

However, LVLMs still suffer from hallucinations, generating responses that describe non-existent objects or incorrectly interpret visual content. Such hallucinations reduce model reliability, motivating extensive research on hallucination mitigation.

Recent studies assume that hallucination-related information can be distinguished in hidden representation spaces and have achieved promising results through contrastive learning and representation editing. However, these approaches typically require additional training and often construct hallucination data based on the observation that hallucinations occur more frequently under degraded image quality.

In this project, we adopt a training-free hallucination mitigation framework and further introduce an additional analysis module to investigate hallucination-related representations. 